% À quoi a servi la dette publique ? Attribution de la dette française aux dépenses qu'elle a financées
% Compilation : pdflatex note-dette.tex (deux passes) depuis fiscalite/dette/
% Chiffres, tableau et figures : FIG_NU=1 python dette/dette.py   (depuis fiscalite/)
\documentclass[11pt,a4paper]{article}
% Préambule commun au brief et à la revue de littérature sur la taxe sur la valeur foncière.
% Même composition que le cours (cours/refaire-l-impot.tex) ; figures « nues » produites par
%   FIG_NU=1 python lvt/graphiques_lvt.py
\usepackage[T1]{fontenc}
\usepackage[utf8]{inputenc}
\usepackage{lmodern}
\usepackage[french]{babel}
\usepackage[a4paper,left=2.8cm,right=2.8cm,top=2.6cm,bottom=2.8cm]{geometry}
\usepackage{graphicx}
\usepackage[export]{adjustbox}
\usepackage{booktabs}
\usepackage{tabularx}
\usepackage{array}
\usepackage{xcolor}
\usepackage{titlesec}
\usepackage[labelfont=bf,labelsep=period,font=small,singlelinecheck=false,
            justification=raggedright,skip=6pt]{caption}
\usepackage{fancyhdr}
\usepackage{enumitem}
\usepackage{microtype}
\usepackage{textcomp}
\usepackage{amsmath}
\providecommand{\euro}{\texteuro}
\usepackage[hidelinks]{hyperref}

\definecolor{encre}{HTML}{1A1A19}
\definecolor{bleu}{HTML}{003399}
\definecolor{vert}{HTML}{1F7A4D}
\definecolor{ocre}{HTML}{C79100}
\definecolor{brique}{HTML}{993333}
\definecolor{filet}{HTML}{C8C6BE}

\titleformat{\section}{\normalfont\large\bfseries\color{bleu}}{\thesection.}{0.6em}{}
\titleformat{\subsection}{\normalfont\normalsize\bfseries\color{bleu}}{\thesubsection}{0.6em}{}
\titlespacing*{\section}{0pt}{2.4ex plus 1ex minus .2ex}{1.2ex plus .2ex}
\titlespacing*{\subsection}{0pt}{2ex plus 1ex minus .2ex}{0.8ex plus .2ex}

\pagestyle{fancy}
\fancyhf{}
\renewcommand{\headrulewidth}{0.4pt}
\fancyhead[R]{\small\thepage}

\setlength{\parindent}{0pt}
\setlength{\parskip}{0.62em}
\linespread{1.05}
\renewcommand{\arraystretch}{1.2}
\setlist{itemsep=2pt,topsep=2pt,parsep=0pt,leftmargin=1.4em}

\newenvironment{encadre}[1]
  {\par\medskip\noindent\begin{minipage}{\linewidth}%
   \color{encre}\setlength{\parskip}{0.5em}%
   \hrule height 0.8pt width \linewidth \vspace{6pt}%
   {\footnotesize\bfseries\color{vert}\MakeUppercase{#1}}\par\vspace{-2pt}%
   \small}
  {\par\vspace{4pt}\hrule height 0.4pt width \linewidth\end{minipage}\par\medskip}

\newcommand{\figbloc}[3]{%
  \begin{figure}[htbp]\centering
  \caption{#2}\label{#3}
  \includegraphics[max width=\linewidth,max height=0.40\textheight]{../figures/nu/#1.pdf}
  \IfFileExists{../figures/nu/#1_note.tex}{\par\vspace{4pt}%
    \parbox{\linewidth}{\raggedright\scriptsize\color{encre}\setlength{\parskip}{2pt}%
    \input{../figures/nu/#1_note.tex}}}{}%
  \end{figure}}

\newcommand{\pc}{\,\%}

\fancyhead[L]{\small\itshape À quoi a servi la dette publique ?}

\begin{document}
\renewcommand{\figurename}{Figure}\renewcommand{\tablename}{Tableau}

\thispagestyle{empty}
\vspace*{0.6cm}
{\footnotesize\color{ocre}\bfseries NOTE DE TRAVAIL\quad·\quad OCTOBRE 2026\par}
\vspace{8pt}
{\fontsize{24}{28}\selectfont\bfseries\color{bleu}À quoi a servi la dette publique ?\par}
\vspace{6pt}
{\large Attribuer la dette française aux dépenses qu'elle a financées\par}
\vspace{8pt}
{\small Colin Baget. Comptes des administrations publiques d'Eurostat et de l'Insee (1949-2025) ;
calculs et figures : \texttt{fiscalite/dette/dette.py}.\par}
\vspace{6pt}
{\color{bleu}\hrule height 1.2pt}
\vspace{4pt}

\section*{L'essentiel}

L'État porte 84~\% de la dette publique, mais ce chiffre dit qui a emprunté, pas pour quoi :
l'État finance une partie des dépenses des collectivités et de la sécurité sociale, par des
transferts et surtout par des impôts qu'il leur cède. Il ne reprend pas pour autant leur dette :
la dette sociale est logée à la CADES, qui fait partie de la sécurité sociale.

La méthode de référence (FMI, Commission européenne) montre pourquoi la dette a monté. De 1995 à
2025, le ratio est passé de 57,5~\% à 115,7~\% du PIB, et les déficits primaires (hors intérêts)
cumulés en expliquent 50,5 points ; les intérêts ont été presque exactement compensés par la
croissance du PIB. Avant 1995, c'était l'inverse : la dette montait surtout parce que les taux
dépassaient la croissance.

Cette méthode ne dit pas quelles dépenses ont fait les déficits. Il faut une convention. Si l'on
tient les recettes pour interchangeables, chaque euro dépensé est emprunté dans la même proportion,
et la dette se répartit au prorata de la dépense. La santé et la protection sociale portent alors
1~960~Md€, soit 57~\% de la dette fin 2025 : 22~\% pour les retraites, 20~\% pour les autres
prestations sociales, 14~\% pour la santé. Ce chiffre reflète surtout la structure de la dépense,
sociale aux trois cinquièmes. Ce qui a changé est plus parlant : depuis 1995, la dépense hors
intérêts a augmenté de 2,3 points de PIB, les recettes de 0,3 point, et la hausse vient entièrement
des retraites (+2,4) et de la santé (+2,0). Ce constat prend les recettes comme référence : rapporté
aux baisses d'impôts décidées, le même déficit s'impute aussi à elles.

Personne n'a publié d'attribution documentée de la dette par fonction ; le rapport Pébereau (2005)
la jugeait impossible. Le chiffre de 57~\% est donc une clé de répartition, pas une mesure de ce que
la dette a causé. Les règles qui confrontent chaque administration à « ses » recettes donnent de
3~\% à 42~\% à la sécurité sociale, selon le statut de la TVA qui lui est affectée.

\section{Ce que dit le chiffre officiel}\label{sec:officiel}

Les comptes nationaux distinguent l'administration centrale (l'État et les organismes divers
d'administration centrale, ODAC), les administrations locales (APUL) et les administrations de
sécurité sociale (ASSO : régimes de base et complémentaires, assurance chômage, CADES, hôpitaux).
Fin 2025, l'administration centrale porte 84~\% de la dette, les collectivités 8~\% et la sécurité
sociale 9~\%, et l'État fait presque tout le déficit (figure~\ref{fig:emetteur}).

\figbloc{dette_emetteur}{Dette de Maastricht par émetteur et solde de chaque sous-secteur,
1995-2025}{fig:emetteur}

L'État finance les autres administrations par trois canaux (figure~\ref{fig:financement}). Les
transferts (dotations, subventions d'équilibre, compensations d'exonérations ciblées) sont des
dépenses de l'État : 55~Md€ vers les collectivités et 23~Md€ vers la sécurité sociale en 2025. Les
impôts qu'il leur cède sont enregistrés comme leurs impôts propres : en 2024, la sécurité sociale a
reçu 58~Md€ de TVA et 63~Md€ d'autres impôts sur la production, les collectivités 52~Md€ de TVA.
La CSG (163~Md€), créée par l'État en 1991, est affectée à la sécurité sociale. Grâce à ces
recettes, sécurité sociale et collectivités paraissent presque à l'équilibre.

\figbloc{dette_financement}{Transferts de l'État et impôts affectés aux autres administrations,
1995-2025}{fig:financement}

Cet enregistrement est une convention discutable : le système européen de comptes (SEC 2010,
§1.78) attribue un impôt à l'administration qui en fixe le taux, et c'est le Parlement qui fixe la
fraction de TVA affectée. L'affectation change aussi souvent. Depuis 2006, l'État compense par
des impôts affectés les allègements de cotisations qu'il décide ; en 2019, il a remplacé le CICE par
une baisse des cotisations patronales compensée par de la TVA ; en 2021 et 2023, la TVA a remplacé
la taxe d'habitation et la CVAE des collectivités.

L'État ne reprend pas la dette des autres. Les 396~Md€ de dette sociale transférés depuis 1996
l'ont été à la CADES, qui les rembourse avec la CRDS et une part de la CSG ; ce remboursement
(environ 16~Md€ par an) n'est pas une dépense en comptabilité nationale, ce qui rend la sécurité
sociale souvent excédentaire. La seule grosse reprise par l'État, 16,8~Md€ en 1993, lui a été
remboursée par la CADES entre 1996 et 2005.

\section{Pourquoi la dette a monté}\label{sec:dynamique}

La décomposition du FMI et de la Commission européenne suit le ratio de dette $d$ (en \% du PIB) :
\[ d_t-d_{t-1} \;=\; \underbrace{pd_t}_{\text{déficit primaire}} \;+\;
\underbrace{\frac{i_t-g_t}{1+g_t}\,d_{t-1}}_{\text{effet « boule de neige »}} \;+\;
\underbrace{sf_t}_{\text{ajustements stock-flux}} \]
où $i_t$ est le taux d'intérêt apparent et $g_t$ la croissance du PIB nominal. La dette augmente
quand on emprunte pour financer un déficit hors intérêts, et quand on emprunte pour payer les
intérêts ; la croissance du PIB, elle, réduit son poids. Le dernier terme regroupe ce qui modifie
la dette sans passer par le déficit (achats d'actifs, trésorerie). Les séries de l'Insee en base
2014 permettent de remonter à 1978, première année où la dette de Maastricht est publiée.

\figbloc{dette_dynamique}{Décomposition de la hausse du ratio de dette, 1978-2025}{fig:dynamique}

Les deux périodes obéissent à des logiques opposées (figure~\ref{fig:dynamique}). De 1978 à 1995,
le ratio passe de 21~\% à 56~\% du PIB, mais les déficits primaires n'y contribuent que pour 8,8
points : après la désinflation, les taux réels étaient élevés, les intérêts (+41,9 points)
dépassaient l'effet de la croissance ($-33{,}4$), et la dette grossissait d'elle-même. Les
ajustements stock-flux pèsent 17,6 points, surtout de 1979 à 1984, sans que ces séries permettent de
les expliquer (opérations financières ou imperfections de la rétropolation, non vérifié).

De 1995 à 2025, le ratio passe de 57,5~\% à 115,7~\% du PIB. Les intérêts (+73,9 points) et la
croissance ($-72{,}3$) se compensent presque : la hausse vient des déficits primaires cumulés (50,5
points) et, pour 6 points, des ajustements stock-flux. Seules les années 1998 à 2001 ont connu un
excédent primaire. Les crises de 2009-2010 et 2020-2021 font 22 des 50,5 points ; le reste vient de
déficits modérés mais presque permanents. La décomposition fait consensus, mais elle s'arrête là : ni le FMI, ni la
Commission, ni les travaux académiques ne répartissent le déficit primaire entre les dépenses.

\section{À quoi a servi la dette : des dépenses fongibles}\label{sec:fongible}

\paragraph{Le principe.} Si aucune recette n'est réservée à une dépense, le déficit est ce qui
manque à une caisse commune, et un euro de retraite n'est pas « plus emprunté » qu'un euro de
gendarmerie : si la dépense fait 100 et le déficit 10, chaque dépense est financée à 10~\% par
l'emprunt. De même, si les collectivités font 20~\% de la dépense, l'État 30~\% et la sécurité
sociale 50~\%, un déficit de 5 points de PIB leur revient pour 1, 1,5 et 2,5 points. L'hypothèse se défend parce que le partage des recettes entre administrations est un
choix du Parlement, révisé presque chaque année. On répartit donc chaque année la hausse de la dette
hors intérêts entre fonctions (classification COFOG, qui range chaque dépense selon son objet,
quel que soit le payeur), au prorata de leur dépense hors intérêts. Les intérêts sont retirés de la
dépense, parce qu'ils sont la conséquence de la dette passée et non une dépense à financer comme les
autres : chaque fonction porte ensuite les intérêts de la dette qui lui revient, au taux moyen payé
par les administrations. Le stock de 1995 est réparti selon la dépense de 1995, ce qui revient à
supposer que la fongibilité valait déjà avant ; comme la santé et la protection sociale faisaient
déjà 55~\% de la dépense en 1995, ce choix pèse peu. Les pensions des fonctionnaires de l'État sont
classées en vieillesse, le RSA versé par les départements en protection sociale.

\figbloc{dette_fonctions}{Dette attribuée à chaque fonction fin 2025, si toutes les dépenses sont
fongibles}{fig:fonctions}

\paragraph{Le résultat.} La santé et la protection sociale portent 1~960~Md€, soit 57~\% de la
dette (figure~\ref{fig:fonctions}) : 766~Md€ pour la vieillesse, 705~Md€ pour les autres
prestations sociales, 488~Md€ pour la santé. Suivent les affaires économiques (12~\%),
l'enseignement (10~\%), les services généraux (9~\%), la défense (4~\%), la police et la justice
(3~\%). Sur la seule dette créée depuis 1995, la part sociale est de 59~\%. Répartie entre
administrations (en retirant de la dépense de chacune les transferts versés aux autres), la dette
revient pour 45~\% à la sécurité sociale, 18~\% aux collectivités et 37~\% à l'État ; l'écart avec
57~\% tient aux dépenses sociales de l'État et des collectivités (pensions de leurs fonctionnaires,
RSA, aides au logement).

\paragraph{La limite.} Avec des dépenses fongibles, chaque fonction porte à peu près sa part de la
dépense : la vieillesse fait 24~\% de la dépense de 2024 et 22~\% de la dette, parce qu'elle
pesait moins dans les années 1990. Le chiffre de 57~\% est donc solide, mais il dit surtout que la
dépense publique est sociale aux trois cinquièmes, pas quelles dépenses ont creusé la dette.

\section{Ce qui a fait monter la dépense}\label{sec:derive}

\figbloc{dette_fonctions_evolution}{Variation de la dépense publique par fonction entre 1995 et
2024}{fig:evolution}

Entre 1995 et 2024, les recettes sont restées à 51~\% du PIB et la dépense hors intérêts a augmenté
de 2,3 points (figure~\ref{fig:evolution}). La vieillesse (+2,4) et la santé (+2,0) ont augmenté de
4,4 points ; toutes les autres fonctions ont baissé ensemble de 2,1 points, affaires économiques,
autres prestations sociales, enseignement et défense en tête. Fipeco trouve des ordres de grandeur
proches sur 1995-2021 (retraites +2,6, santé +1,3) ; l'écart sur la santé tient aux bornes, la
santé ayant gagné 0,8 point depuis 2019.

La méthode d'Auerbach (1994) traduit cette dérive en dette. Le déficit primaire de chaque année
s'écrit comme celui d'une année de référence, plus la hausse de chaque dépense, moins la hausse des
recettes ; cumulé sur 1996-2025, il ventile exactement les 50,5 points de déficits primaires. Avec
1997 comme référence, année où le solde primaire était presque à l'équilibre, la hausse de la
vieillesse et de la santé représente 74 points cumulés, plus que la totalité des déficits
primaires, et les autres dépenses ont baissé de 23 points (figure~\ref{fig:derive}). Le partage
dépend de l'année de référence : prendre 1995 (déficit primaire de 1,6 point) ou 2000 (excédent de
1,6 point) change le signe des autres postes. Seules la vieillesse (+42 à +55 points) et la santé
(+32) pèsent lourd quelle que soit la référence.

\figbloc{dette_derive}{Dérive de chaque poste par rapport à 1997, cumulée depuis 1996}{fig:derive}

Ce constat prend les recettes comme référence, ce qui est un choix. Champsaur et Cotis (2010)
chiffrent à environ 3 points de PIB les baisses de prélèvements décidées entre 1999 et 2008, sans
lesquelles la dette aurait été plus faible d'environ 20 points ; ils concluent que la dégradation
« est tout autant due à une maîtrise insuffisante de la dépense qu'à des allègements de
prélèvements ». Depuis 2019, les recettes sont passées de 53,0~\% à 51,2~\% du PIB. Imputer la dérive
à la dépense sociale ou aux baisses d'impôts sont deux lectures également comptables du même
déficit, qui diffèrent par la référence. Aucune ne dit non plus ce qu'il aurait fallu faire : le
vieillissement augmente mécaniquement la dépense de retraite et de santé, et le choix entre pensions,
prélèvements et autres dépenses est politique.

Avant 1995, l'Insee ne publie pas de dépense par fonction, dans aucune de ses bases ; on peut
seulement suivre les prestations sociales (en espèces, et remboursements de soins et aides au
logement). Elles passent de 11,3~\% du PIB en 1960 à 17,9~\% en 1978 et 25,8~\% en 2025
(figure~\ref{fig:prestations}). Depuis 1978, le déficit primaire n'a augmenté que de 2,2 points :
l'essentiel de la hausse des prestations a été couvert par des hausses de prélèvements ou des
baisses d'autres dépenses.

\figbloc{dette_prestations}{Prestations sociales et déficit primaire des administrations,
1960-2025}{fig:prestations}

\section{Si l'on refuse la fongibilité}\label{sec:variantes}

On peut aussi demander de combien la dette serait plus faible si la sécurité sociale avait dépensé
moins, en tenant pour acquis que l'État lui aurait cédé les mêmes impôts. Cette question conduit à
des règles qui confrontent chaque administration à ses propres recettes
(tableau~\ref{tab:attribution}). Dans la plus simple (M2), chaque administration garde son déficit,
et la part qui finance des transferts suit l'argent : si l'État emprunte 10 et consacre 20~\% de ses
dépenses à la sécurité sociale, 2 lui sont attribués. La sécurité sociale, presque à l'équilibre
depuis trente ans, ne porte alors que 3~\% de la dette créée depuis 1995. Si l'on traite la TVA
affectée comme un transfert de l'État (M3), sa part monte à 42~\%, et elle dépasse 100~\% si l'on
traite de même tous ses impôts indirects (M3+). La CSG reste dans tous les cas une ressource
propre ; la traiter comme un impôt d'État attribuerait à la sécurité sociale plus que la totalité
du déficit.

\begin{table}[htbp]\centering\small
\caption{Part de chaque administration (ou du social) dans la dette fin 2025, selon la règle (en \%)}
\label{tab:attribution}
\setlength{\tabcolsep}{4pt}
\begin{tabular}{@{}l>{\raggedright\arraybackslash}p{5.6cm}rrrrrr@{}}
\toprule
& & \multicolumn{3}{c}{Dette totale} & \multicolumn{3}{c}{Dette créée depuis 1995} \\
\cmidrule(lr){3-5}\cmidrule(l){6-8}
& Règle & ASSO & APUL & État & ASSO & APUL & État \\
\midrule
M0 & Émetteur & 9 & 9 & 82 & $-$1 & $-$6 & 106 \\
M2 & Transferts explicites suivis jusqu'au dépensier & 10 & 14 & 76 & 3 & 6 & 91 \\
M3 & Idem, TVA affectée traitée comme un transfert de l'État & 26 & 19 & 55 & 42 & 19 & 39 \\
M3+ & Idem, tous impôts indirects affectés à la sécurité sociale & 61 & 15 & 25 & 131 & 8 & $-$39 \\
M1 & Prorata de la dépense primaire finale & 27 & 19 & 54 & 46 & 20 & 34 \\
M4 & Prorata de la dépense par fonction (santé et protection sociale) & 32 &  & 68 & 59 &  & 41 \\
\bottomrule

\end{tabular}
\par\vspace{4pt}
\parbox{\linewidth}{\raggedright\scriptsize\color{encre} Source : Eurostat, Insee ; calculs de
l'auteur (\texttt{dette/dette.py}). ASSO : sécurité sociale ; APUL : collectivités ; État : État
et ODAC. La dette créée depuis 1995 est la somme des déficits primaires de 1996 à 2025 et des
intérêts qu'ils ont générés. Dans la dette totale, le stock de 1995 et les ajustements stock-flux
sont répartis comme la dépense pour les règles fongibles (M1, M4) et laissés à l'émetteur pour les
autres. Pour M4, la première colonne est la santé et la protection sociale, la troisième les autres
fonctions. Une part négative signale des excédents primaires cumulés.}
\end{table}

Tout se joue sur le statut des impôts affectés, et c'est ce qui rend ces règles fragiles. La
transformation du CICE en 2019 a remplacé une dépense de l'État par une baisse de cotisations
compensée par environ 35~Md€ de TVA par an, sans changer les prestations ; dans la règle M3, elle
fait pourtant passer 35~Md€ de déficit par an de l'État à la sécurité sociale. Une règle cohérente
compterait la TVA affectée comme une ressource de l'État et, symétriquement, les cotisations
exonérées comme une ressource de la sécurité sociale ; les deux corrections se compenseraient pour
l'essentiel, et le résultat serait proche de M2. Selon la Cour des comptes, la compensation par la
TVA a même été défavorable à la sécurité sociale, de 18~Md€ cumulés entre 2019 et 2024 (cité par
Bach et al., 2026). Ces règles répondent à une question de
gouvernance (quelle administration est en déficit avec les recettes que la loi lui donne ?), pas à
celle de l'usage de la dette.

\section{Ce qu'en disent les travaux existants}\label{sec:litterature}

La recherche menée pour cette note (Cour des comptes, Insee, France Stratégie, OFCE, Fipeco,
travaux américains, britanniques et allemands, littérature sur la dynamique de la dette) n'a trouvé
aucune attribution documentée du stock de dette à des fonctions. Le rapport Pébereau (2005) la
jugeait impossible : « il n'existe qu'un seul déficit […], sans que l'on sache identifier précisément
les dépenses qui ont été financées par de la dette ». Le rapport n'admettait que deux exceptions, la
sécurité sociale et l'assurance chômage, dont la dette est identifiée, et ne ventilait la dette que
selon sa contrepartie : en moyenne, 40~\% des déficits de 1978 à 2002 n'avaient financé aucun actif.
Notre prorata respecte le principe d'un seul déficit, mais c'est une clé de répartition, pas une
identification.

Les travaux établis décomposent des variations (Fipeco ; Auerbach, 1994 ; Office for Budget
Responsibility, 2024) ou chiffrent des contrefactuels, qui donnent des résultats opposés selon le
poste qu'on fait varier. Champsaur et Cotis (2010) raisonnent côté recettes (environ 20 points de
dette en moins sans les baisses de prélèvements depuis 1999). Sylvain Catherine raisonne côté
dépenses : avec la dépense de retraite allemande depuis 1995, la dette serait de 57,5~\% du PIB en
2025 au lieu de 115,7~\% (calcul publié sur X, connu par sa reprise dans \emph{JeChange}). Les deux
calculs sont corrects, mais ces effets ne s'additionnent pas à 100~\%. Aux États-Unis, les mêmes
projections du Congressional Budget Office conduisent Riedl (2020) à imputer 97~\% des déficits
futurs aux retraites et à Medicare, et Kogan (2023) à imputer 57~\% de la hausse de la dette aux
baisses d'impôts : le résultat dépend de la référence.

Jean-Pascal Beaufret (2023) estime que les retraites « contribuent à plus de la moitié des déficits
publics ». Il compte comme déficit toute la dépense de pensions non couverte par des cotisations au
taux de droit commun ou des impôts affectés (71~Md€ en 2021, dont 45~Md€ de surcotisations des
employeurs publics) ; le Haut-Commissariat au Plan (2022) trouvait environ 30~Md€ avec une méthode
voisine. Le calcul montre que l'équilibre des régimes de fonctionnaires repose sur des
taux de cotisation employeur de 74~\% (civils de l'État) ou 126~\% (militaires), mais il refuse la
fongibilité : appliqué à la défense, qui n'a aucune recette affectée, il lui attribuerait toute sa
dépense. Le Conseil d'orientation des retraites répond qu'un solde n'est univoque qu'entre des
« frontières étanches », et qu'il y a sinon « autant de conventions […] que d'hypothèses ». Sylvain
Catherine défend une position proche : en avril 2026, il a écrit que les retraites étaient « la
cause […] de la moitié de l'endettement public », ce qu'Henri Sterdyniak a contesté. Nicolas Dufourcq (2025) avance que « 2~000 milliards au minimum » de la
dette sont des prestations sociales versées à crédit, sans exposer sa méthode, et ses propres chiffres
(888~Md€ de prestations sur plus de 1~600~Md€ de dépense) ne concordent pas avec les « 60~\% de la
dette » qu'il avance en entretien ; notre calcul explicite donne 1~960~Md€. Enfin, l'Institut des politiques publiques (Bach et al., 2026) et
François Ecalle (Fipeco) jugent que les soldes par administration n'ont guère de sens : attribuer un
solde à une administration suppose qu'elle maîtrise ses dépenses et ses recettes, ce qui n'est le
cas ni de la sécurité sociale ni des collectivités. Seule compte la dette de l'ensemble, ce qui
rejoint le point de départ de cette note.

\section{Limites}\label{sec:limites}

La dépense par fonction n'existe qu'à partir de 1995, chez l'Insee comme chez Eurostat, et s'arrête
en 2024 (2025 reprend la structure de 2024) ; le stock de 1995 est réparti selon la dépense de 1995.
Pour les fonctionnaires de l'État, la comptabilité nationale compte deux fois la charge des
retraites, en pensions (vieillesse) et en cotisation employeur fictive (fonction de l'employeur) ;
corriger cette convention retirerait 28,9~Md€ de dépense en 2023, surtout en enseignement, défense
et ordre public (CAE, 2025), et relèverait d'environ un point la part sociale. La dette est brute :
les réserves des régimes complémentaires ne sont pas déduites.

\appendix
\section{Méthode}\label{sec:annexe}

\paragraph{Données.} Eurostat : dépense par fonction COFOG (\texttt{gov\_10a\_exp}, 1995-2024),
comptes des administrations par sous-secteur, dont les transferts versés à chaque autre sous-secteur
(\texttt{gov\_10a\_main}), impôts reçus (\texttt{gov\_10a\_taxag}), PIB. Insee : dette de
Maastricht par sous-secteur (contributions consolidées, base 2020) ; comptes nationaux en base 2014
(séries arrêtées) pour la dette (depuis 1978), le déficit, les intérêts, les prestations sociales
et le PIB (depuis 1949). Chaque variation est calculée dans une seule base : base 2014 jusqu'en 1995,
base 2020 ensuite ; le changement de base relève la dette de 1995 de 1,4 point.

\paragraph{Dérive.} En notant $x_{jt}$ la dépense primaire de la fonction $j$ et $\tau_t$ les
recettes, en \% du PIB, et $R$ l'année de référence, $pd_t=pd_R+\sum_j(x_{jt}-x_{jR})-(\tau_t-\tau_R)$.

\paragraph{Méthode fongible (M4, M1).} Soit $D_t$ la dette, $I_t$ les intérêts, $r_t=I_t/D_{t-1}$
et $s_{jt}$ la part de la fonction $j$ dans la dépense hors intérêts. La dette attribuée suit
$D_{jt}=(1+r_t)D_{j,t-1}+s_{jt}(\Delta D_t-I_t)$, avec $D_{j,1995}=s_{j,1995}D_{1995}$ ; la somme
des $D_{jt}$ est $D_t$. M1 remplace la fonction par le sous-secteur, en retirant de sa dépense les
transferts versés aux autres.

\paragraph{Règles non fongibles (M0, M2, M3, M3+).} M0 attribue à chaque sous-secteur son déficit
primaire. Dans M2, soit $d_j$ le déficit primaire du sous-secteur $j$, $\tau_{ij}$ la part de ses
emplois que $i$ transfère à $j$ et $\sigma_j=1-\sum_k\tau_{jk}$ ; l'emprunt dont dispose $j$
vérifie $e_j=d_j+\sum_i e_i\tau_{ij}$, et on lui attribue $\sigma_je_j$. Dans M3 et M3+, l'impôt
affecté est retiré des recettes du bénéficiaire et versé par l'État comme un transfert. Le stock de
1995 et les ajustements stock-flux restent à l'émetteur.

\section*{Références}
\begin{small}\setlength{\parskip}{3pt}\raggedright
Auerbach, A. (1994), « The U.S. Fiscal Problem: Where We Are, How We Got Here, and Where
We're Going », \emph{NBER Macroeconomics Annual}, vol.~9, p.~141-186.

Bach, L., Bozio, A., De~Brouwer, O. et Eryzhenskiy, I. (2026), « Qui est responsable du déficit
public ? », in \emph{Perspectives budgétaires}, Institut des politiques publiques, juin 2026, chap.~1.

Beaufret, J.-P. (2023), « Retraites obligatoires et déficits publics. Pour la clarté »,
\emph{Commentaire}, n°~182, p.~245-254.

Catherine, S. (2026), calcul publié sur X, repris par F. Tostivint, \emph{JeChange}, 18~août
2026 (source secondaire).

Champsaur, P. et Cotis, J.-P. (2010), \emph{Rapport sur la situation des finances publiques},
avril 2010, La Documentation française.

Commission sur la dette publique, présidée par M. Pébereau (2005), \emph{Rompre avec la
facilité de la dette publique}, La Documentation française.

Conseil d'analyse économique (2025), \emph{Focus} n°~121.

Conseil d'orientation des retraites (2023), séance du 21 septembre 2023, documents de travail.

Dufourcq, N. (2025), \emph{La dette sociale de la France 1974-2024}, Odile Jacob (non lu ; cité
par la presse).

Ecalle, F. (2023), « Cinquante ans de dépenses publiques en France (1972-2022) », Fipeco,
2~novembre 2023.

Eurostat (2013), \emph{Système européen des comptes, SEC 2010}, §1.78.

Kogan, B. (2023), « Tax Cuts Are Primarily Responsible for the Increasing Debt Ratio », Center
for American Progress, mars 2023.

Office for Budget Responsibility (2024), \emph{Fiscal risks and sustainability}, septembre 2024.

Riedl, B. (2020), « Behind CBO's \$100 Trillion in Projected Deficits over 30 Years », Manhattan
Institute, novembre 2020.
\end{small}

\end{document}
